\subsection{From Classification to Quantification: The Dataset Shift Challenge}
\label{sec:5.5}

The previous section surveyed prior efforts to identify dark matter subhalos among Fermi-LAT unassociated sources, all of which relied on a common strategy: classify individual sources and count the candidates. We now examine why this approach is fundamentally limited and introduce the alternative framework---quantification learning---that underpins the analysis presented in the remainder of this chapter. The formal vocabulary of dataset shift was introduced in Section~\ref{sec:3.4}; this section applies that framework to the specific problem of subhalo searches and motivates the construction of a generative mixture model capable of simultaneously fitting class prevalences, correcting for observational selection effects, and constraining the dark matter annihilation cross-section from the data.


\subsubsection{Why Standard Classification Fails for DM Subhalo Searches}
\label{sec:5.5.1}

The classify-and-count paradigm described in Section~\ref{sec:5.4} rests on a conceptually simple idea: train a classifier to distinguish dark matter subhalos from astrophysical sources, apply it to the unassociated catalog, count the resulting candidates, and derive upper limits on $\langle\sigma v\rangle$ under the assumption that the true number of DM subhalos cannot exceed the number of candidate sources~\cite{2019JCAP...07..020C, 2023JCAP...07..033B}. In practice, however, this strategy suffers from three interrelated problems that compromise the statistical rigor of the resulting constraints.

The first and most fundamental difficulty is the absence of verified training examples for the signal class. No dark matter subhalo has ever been identified in gamma-ray data, so the ``DM'' class used in supervised classification has no empirical anchor. Researchers instead generate synthetic DM spectra from theoretical models and inject them into the training set as surrogates for real subhalos~\cite{2023JCAP...07..033B, 2023MNRAS.520.1348G}. While this approach provides the classifier with some notion of how a DM source might appear in feature space, the predicted probabilities it assigns to individual sources cannot be calibrated against reality. Calibration requires a testing set whose true labels are known, and for a class with zero confirmed members, no such set exists~\cite{Amerio:2025fhz}. The classifier may learn the broad features of simulated spectra, but without calibration, the probability that a given source is a DM subhalo has no well-defined statistical meaning.

The second issue concerns the treatment of class prevalences during training. Standard supervised learning algorithms perform best when the training classes are balanced---roughly equal numbers of examples per class---and practitioners accordingly construct training sets with comparable fractions of AGNs, pulsars, and simulated DM sources, typically around one-third each~\cite{Amerio:2025fhz}. This artificial equalization bears no relation to the actual composition of the gamma-ray sky. Among associated sources at $|b| > 10^\circ$, AGNs outnumber pulsars by nearly an order of magnitude, and the DM contribution, if present, is far smaller still~\cite{Amerio:2025fhz}. A classifier trained on balanced classes will systematically overpredict the minority class. As a concrete illustration, consider the simpler task of classifying sources into pulsars and AGNs: if one balances the training set at 50/50, the classifier predicts a pulsar fraction among unassociated sources that is substantially larger than what the known source population would suggest~\cite{Amerio:2025fhz}. The same inflation is expected---and indeed far more severe---for the DM class, whose true prevalence may be zero. In the language of dataset shift (Section~\ref{sec:3.4}), this constitutes a prior shift: the class prevalences $p(k)$ differ between the training set and the target population, and the classifier has been given no mechanism to correct for this mismatch.

The third weakness is the arbitrary threshold by which candidates are selected. In machine learning classification, label assignment requires converting a continuous probability score into a discrete decision: a source is labeled a DM candidate if $p_\mathrm{DM}(\mathbf{x}) > p_\mathrm{threshold}$, where $p_\mathrm{threshold}$ is chosen by the analyst. No statistical principle determines this threshold. A lenient choice (e.g., $p_\mathrm{DM} > 0.5$) produces a large candidate list and correspondingly weak $\langle\sigma v\rangle$ bounds, because the analysis must accommodate the possibility that many bright subhalos exist. A stringent choice (e.g., $p_\mathrm{DM} > 0.9$) yields few candidates and tighter bounds, but excludes potentially genuine signals. The resulting physical constraints therefore depend on a free parameter that carries no physical or statistical justification~\cite{Amerio:2025fhz}.

These problems compound in practice. Butter et al.~\cite{2023JCAP...07..033B} applied Bayesian neural networks to the 4FGL-DR4 catalog and identified 281 DM subhalo candidates in their most conservative scenario---a number that directly reflects the combined effect of balanced training and a particular probability threshold. Different choices would have yielded a different candidate count and, consequently, different $\langle\sigma v\rangle$ bounds~\cite{Amerio:2025fhz}. Yet no mechanism exists within the classify-and-count framework to determine which threshold, if any, produces the correct answer. The upper limits derived in this way lack a well-defined confidence level in the frequentist sense: they represent an assumption ($N_\mathrm{DM} \leq N_\mathrm{candidates}$) rather than a statistical inference.

These limitations are not unique to machine-learning-based searches. The hand-crafted approaches of Coronado-Bl\'azquez et al.~\cite{2019JCAP...07..020C, 2019JCAP...11..045C} also adopt a classify-and-count logic, selecting candidate subhalos on the basis of spectral, spatial, and variability criteria and implicitly assuming that the training sample (associated sources with known properties) represents the full population. While the individual selection criteria differ, the structural problem is the same: the analysis cannot separate genuine DM contributions from the statistical effects of prior and covariate shifts, and the resulting bounds depend on analyst-chosen thresholds rather than on a coherent statistical model of the data.

The failure of classify-and-count is therefore not a technical limitation of any particular classifier, but a conceptual one: the approach attempts to answer a population-level question---how many DM subhalos contribute to the unassociated catalog?---using a tool designed for individual-level decisions. A different framework is needed, one that models the composition of the unassociated population as a whole and fits the DM contribution directly to the data. Before introducing that framework, we first examine in detail the dataset shift between associated and unassociated sources that any such model must address.


\subsubsection{Dataset Shift in the Fermi-LAT Context}
\label{sec:5.5.2}

The shortcomings described in the preceding section are not merely theoretical: the associated and unassociated Fermi-LAT source populations are measurably different, and these differences violate the assumptions on which standard classifiers rely. We introduced the formal dataset shift framework in Section~\ref{sec:3.4}; here, we examine how prior and covariate shifts manifest concretely in the 4FGL-DR4 catalog, and why their combined effect presents the central statistical challenge of the subhalo search.

Each source in the 4FGL-DR4 catalog is characterized by a log-parabola energy spectrum with three parameters: the integrated flux $\log_{10}\phi$, the spectral index $\alpha$, and the spectral curvature $\beta$~\cite{2023arXiv230712546B}. The distributions of these parameters differ visibly between associated and unassociated sources (see Figure~\ref{fig:assoc_unas_distributions} in Section~\ref{sec:3.4}). This mismatch has two possible origins---a change in the relative composition of source classes (prior shift), a change in the feature distributions within each class due to selection effects (covariate shift)---or some combination of both~\cite{MorenoTorres2012AUV, Amerio:2025fhz}.

Prior shift, introduced formally in Section~\ref{sec:3.4.2} (Equation~\ref{eq:prior_shift}), corresponds to a change in class prevalences $p(k)$ between the training and target populations, while the class-conditional feature distributions $p(\mathbf{x}|k)$ remain the same. In the Fermi-LAT context, this effect is expected on physical grounds. The association procedure links gamma-ray sources to multi-wavelength counterparts, but this linkage is far from uniform across source classes. Bright blazars at high Galactic latitudes are easily identified through optical and radio surveys, so they are efficiently associated. Pulsars, by contrast, require dedicated radio timing campaigns or gamma-ray periodicity searches, and many pulsars with unfavorable beaming geometries remain undetected at other wavelengths~\cite{2023RASTI...2..735M}. As a result, Galactic sources constitute roughly 29\% of the unassociated population at $|b| > 10^\circ$ but only about 6\% of the associated sample~\cite{Amerio:2025fhz}. This compositional difference is a textbook example of prior shift: the intrinsic spectral properties of pulsars and AGNs have not changed, but their relative weights in the two samples are different.

Covariate shift (Section~\ref{sec:3.4.2}, Equation~\ref{eq:covariate_shift}) is the complementary effect: the class prevalences $p(k)$ are the same, but the feature distributions $p(\mathbf{x})$ differ because of selection biases. In the Fermi-LAT catalog, association is strongly correlated with source brightness. Bright sources have small localization error regions, making multi-wavelength counterpart identification straightforward; faint sources, conversely, are often left unassociated regardless of their astrophysical nature. As a result, the unassociated population is systematically shifted toward lower fluxes and larger spectral uncertainties compared to the associated sample~\cite{2023RASTI...2..735M, Amerio:2025fhz}. The spectral index and curvature distributions are also affected: the ratio of unassociated to associated source densities varies monotonically across portions of the $(\alpha, \beta)$ feature space, reflecting the flux-dependent completeness of the association procedure~\cite{Amerio:2025fhz}.

The difficulty is that these two effects are entangled. A change in the marginal class prevalences $p(k)$ necessarily produces a corresponding change in the marginal feature distribution $p(\mathbf{x})$, since different source classes occupy different regions of feature space. Conversely, a selection bias that alters $p(\mathbf{x})$ can mimic a change in class fractions when projected onto summary statistics. The same observed mismatch between associated and unassociated distributions can therefore be explained by prior shift, by covariate shift, or by a combination of both~\cite{MorenoTorres2012AUV, Amerio:2025fhz}---a degeneracy already discussed in general terms in Section~\ref{sec:3.4.2}. In the context of the subhalo search, this degeneracy acquires a third dimension: the discrepancy could also arise, in part, from a genuinely new source class---dark matter subhalos---contributing to the unassociated population with zero representation among associated sources. A DM subhalo population constitutes a particular case of prior shift: the DM prevalence goes from exactly zero in the training set to some unknown, possibly non-zero value in the target population~\cite{Amerio:2025fhz}.

Ignoring this three-way degeneracy is dangerous regardless of which component is underestimated. A classifier that assumes $p_\mathrm{train}(\mathbf{x}, k) = p_\mathrm{target}(\mathbf{x}, k)$ will misinterpret the excess of pulsar-like unassociated sources: it may attribute them entirely to pulsars (underestimating a possible DM contribution) or, if the DM class template overlaps with the pulsar-like region of feature space, partly to DM subhalos (overestimating the DM contribution). Neither outcome is controlled by the data---both reflect the implicit prior assumptions baked into the training procedure. Breaking the degeneracy requires a model flexible enough to fit prior shift, covariate shift, and a potential new-class contribution simultaneously, with all three effects determined by the data rather than assumed in advance~\cite{Amerio:2025fhz}. We describe such a model in the following sections.


\subsubsection{Quantification Learning: From $p(k|\mathbf{x})$ to $p(\mathbf{x}|k)$}
\label{sec:5.5.3}

The shortcomings of classify-and-count trace back to a single conceptual choice: standard classifiers estimate the posterior probability $p(k|\mathbf{x})$ of a source belonging to class $k$ given its observed features $\mathbf{x}$. This discriminative approach is effective for labeling individual objects, but it does not provide a coherent model of the population. In particular, it offers no natural mechanism for estimating class prevalences, correcting for dataset shift, or constructing a likelihood that can be maximized to constrain physical parameters such as $\langle\sigma v\rangle$~\cite{Amerio:2025fhz}.

An alternative is to work in the generative direction: instead of asking ``what class does this source belong to?'', one asks ``what is the probability of observing features $\mathbf{x}$ given that the source belongs to class $k$?'' This amounts to estimating the class-conditional densities $p(\mathbf{x}|k)$~\cite{Amerio:2025fhz}. The two quantities are related through Bayes' theorem,
\begin{equation}
p(k|\mathbf{x}) = \frac{p(\mathbf{x}|k)\,p(k)}{p(\mathbf{x})},
\label{eq:bayes_generative}
\end{equation}
so the generative approach contains strictly more information: from $p(\mathbf{x}|k)$ and the class priors $p(k)$, one can recover $p(k|\mathbf{x})$, but the reverse is not generally possible without additional assumptions~\cite{Amerio:2025fhz}.

This shift in perspective is at the core of \textit{quantification learning}, a field in machine learning concerned with estimating class prevalences in unlabeled datasets rather than classifying individual instances~\cite{10.1145/3117807_Gonzalez_quantification, esuli2023learning}. The central insight is direct: if one can construct the class-conditional densities $p(\mathbf{x}|k)$ from a labeled training set, then the distribution of the unlabeled target data can be written as a mixture,
\begin{equation}
p_\mathrm{target}(\mathbf{x}) = \sum_k \pi_k\, p(\mathbf{x}|k),
\label{eq:mixture_generic}
\end{equation}
where the mixing weights $\pi_k$ represent the unknown class prevalences. These weights are free parameters, fit to the target data by maximizing the model likelihood~\cite{10.1145/3117807_Gonzalez_quantification, 2024arXiv240100490M}. The class prevalences are thereby determined by the data, not assumed during training---exactly the property that classify-and-count methods lack.

In the quantification learning literature, this procedure is known as \textit{template matching}: the class-conditional densities serve as templates, and the mixing weights are adjusted so that the weighted combination best reproduces the observed target distribution~\cite{2024arXiv240100490M}. The concept is familiar in astrophysics. Spatial analyses of diffuse gamma-ray emission, for example, routinely decompose observed maps into a linear combination of template maps tracing distinct physical components---hadronic emission, inverse Compton scattering, isotropic background---and fit the template normalizations to the data. The approach adopted here is the spectral-parameter analogue: the ``templates'' are probability density functions of the spectral parameters $(\log_{10}\phi, \alpha, \beta)$ for each source class, and the ``normalizations'' are the class prevalences $\pi_k$~\cite{Amerio:2025fhz}.

The generative formulation brings three concrete advantages over classification for the subhalo search. First, it yields a well-defined likelihood: the product $\mathcal{L} = \prod_i p_\mathrm{target}(\mathbf{x}_i)$ over all unassociated sources provides a global goodness-of-fit measure that can be maximized to determine the best-fit parameters and profiled to derive confidence intervals on $\langle\sigma v\rangle$~\cite{Amerio:2025fhz}. Second, the model is generative in the statistical sense: once the parameters are fixed, it can be sampled to produce synthetic datasets, enabling Monte Carlo validation of the inference procedure and estimation of statistical uncertainties~\cite{Amerio:2025fhz}. Third, because the class prevalences are fit rather than assumed, the model naturally accommodates prior shift and can include a new class---dark matter subhalos---whose prevalence was zero in the training sample but may be non-zero in the target population~\cite{Amerio:2025fhz, 10.1162/089976602753284446_Saerens}.

The practical construction of $p(\mathbf{x}|k)$ from the associated source catalog, and the extension of the mixture model to simultaneously handle covariate shift and the DM component, are described in the next section.


\subsubsection{The Mixture Model Concept}
\label{sec:5.5.4}

We now outline the statistical model that implements the quantification learning framework for the subhalo search, deferring the detailed construction, optimization procedure, and results to the paper body (Section~\ref{sec:5.6} onward).

The unassociated source population at $|b| > 10^\circ$ is modeled as a three-component mixture. The first two components represent the known astrophysical source classes---Galactic (primarily pulsars, supernova remnants, and binaries) and extragalactic (predominantly AGNs)---whose spectral parameter distributions are estimated from the associated catalog using kernel density estimation~\cite{Amerio:2025fhz, bishop}. The third component represents a hypothetical population of dark matter subhalos, whose spectral parameter distribution is not drawn from data but instead derived from Monte Carlo simulations for fixed DM parameters (mass $m_\mathrm{DM}$, annihilation cross-section $\langle\sigma v\rangle$, and annihilation channel)~\cite{Amerio:2025fhz}. This physically anchored construction means that the DM template is not a free-form density: it is fully determined by the assumed particle physics and the J-factor distributions obtained from N-body simulations (Section~\ref{sec:5.2}).

The model for the probability density of an unassociated source with feature vector $\mathbf{x} = (\log_{10}\phi, \alpha, \beta)$ takes the form introduced in Section~\ref{sec:3.4} (see Equations~\ref{eq:joint_factorizations} and~\ref{eq:sigmoid_modulation}):
\begin{equation}
\tilde{p}_\mathrm{unas}(\mathbf{x}) = \left(\sum_{k=1}^{2} \pi_k\, p_\mathrm{assoc}(\mathbf{x}|k)\right) \tilde{C}(\mathbf{x};\boldsymbol{\theta}_\mathrm{cov}) + \pi_\mathrm{DM}(\boldsymbol{\theta}_\mathrm{DM})\, p_\mathrm{DM}(\mathbf{x};\boldsymbol{\theta}_\mathrm{DM}).
\label{eq:mixture_model}
\end{equation}
Each component in this equation addresses one of the problems identified in the previous sections:

The class prevalences $\pi_k$ for the Galactic and extragalactic components are free parameters, fit to the unassociated data by maximum likelihood. This handles the prior shift directly: the model determines how much of the unassociated population is Galactic and how much is extragalactic, rather than inheriting those fractions from the associated catalog~\cite{10.1145/3117807_Gonzalez_quantification, 10.1162/089976602753284446_Saerens}.

The covariate shift, arising from the brightness-dependent association bias, is corrected by the modulation function $\tilde{C}(\mathbf{x};\boldsymbol{\theta}_\mathrm{cov})$, defined in Section~\ref{sec:3.4.2} (Equation~\ref{eq:sigmoid_modulation}) as a product of sigmoid functions along each feature dimension. The sigmoid parameters $\boldsymbol{\theta}_\mathrm{cov} = \{b_d, c_d\}$ are optimized jointly with the prevalences, so the model adapts the shape of the astrophysical templates to account for the fact that unassociated sources are systematically fainter and have different spectral distributions than their associated counterparts~\cite{Amerio:2025fhz}.

The DM component $p_\mathrm{DM}(\mathbf{x};\boldsymbol{\theta}_\mathrm{DM})$ is constructed through a forward simulation pipeline. For a given $m_\mathrm{DM}$ and $\langle\sigma v\rangle$, J-factors are sampled from the distributions provided by N-body simulations, the corresponding gamma-ray spectra are computed via the annihilation flux formula (Section~\ref{sec:5.3}), and the resulting sources are injected into a realistic Fermi-LAT simulation to determine which subhalos would be detected above the $\mathrm{TS} > 25$ threshold. The surviving sources are fit with log-parabola spectra, producing the DM density in feature space. The DM weight $\pi_\mathrm{DM}$ is not a free parameter in the same sense as the astrophysical prevalences: it is fixed by the physics,
\begin{equation}
\pi_\mathrm{DM}(\boldsymbol{\theta}_\mathrm{DM}) = \frac{N_\mathrm{sub}^\mathrm{ROI}(\boldsymbol{\theta}_\mathrm{DM})}{N_\mathrm{unas}^\mathrm{ROI}},
\label{eq:pi_dm}
\end{equation}
where $N_\mathrm{sub}^\mathrm{ROI}$ is the expected number of detectable subhalos in the region of interest and $N_\mathrm{unas}^\mathrm{ROI}$ is the total number of unassociated sources. The DM contribution therefore increases with $\langle\sigma v\rangle$ and depends on the J-factor model, but it is not adjusted independently of the underlying physics~\cite{Amerio:2025fhz}.

The full likelihood is the product of $\tilde{p}_\mathrm{unas}(\mathbf{x}_i)$ over all unassociated sources in the region of interest. This likelihood is maximized using the Expectation-Maximization algorithm~\cite{10.1162/089976602753284446_Saerens, bishop}, simultaneously optimizing $\pi_k$, $\boldsymbol{\theta}_\mathrm{cov}$, and $\langle\sigma v\rangle$. Upper bounds on the annihilation cross-section are obtained by profiling the likelihood ratio test statistic and identifying the value of $\langle\sigma v\rangle$ at which $\mathrm{TS} = 3.84$, corresponding to the 95\% confidence level for one degree of freedom~\cite{Amerio:2025fhz}.

This framework addresses each of the four shortcomings identified earlier. It eliminates the need for a probability threshold, since the DM contribution is a continuous parameter fit by maximum likelihood rather than a binary classification decision. It handles prior shift by fitting class prevalences and covariate shift by fitting the sigmoid modulation, breaking the degeneracy between the two effects. And it provides statistically well-defined confidence intervals through the profile likelihood, replacing the classify-and-count assumption $N_\mathrm{DM} \leq N_\mathrm{candidates}$ with a rigorous inference of the DM population size from the data~\cite{Amerio:2025fhz}.

The remainder of this chapter presents the full analysis. Section~\ref{sec:5.6} details the statistical model construction, including the data selection, KDE estimation of the astrophysical templates, and the covariate shift parametrization. Section~\ref{sec:5.7} describes the DM subhalo simulation pipeline and J-factor distributions. Section~\ref{sec:5.8} presents the mixture model results and the constraints on the DM annihilation cross-section. Section~\ref{sec:5.9} discusses the implications and places the results in the context of other indirect detection limits.
